% SGA 3, Expose VI_A, Loop 1 English reconstruction.
% Scope: paragraph 5.3 and subparagraphs 5.3.1--5.3.3 in full;
% paragraph 5.4 is excluded.
% Authority: Polo--Gille Exp6A-13oct24.pdf, component-local pages 31--32
% (combined-reader pages 334--335; printed pages 323--324).
% OCR and the Reinhold English Markdown are comparison witnesses only.

\medskip
\phantomsection\label{par:sga3-VIA-5.3}%
\noindent\textbf{5.3.---}
现在回顾第~IV~篇的若干陈述.

\medskip
\phantomsection\label{par:sga3-VIA-5.3.1}%
\noindent\textbf{5.3.1.---}
\SGARefLink{sga3:iv:proposition:5:2:7}{第~IV~篇 5.2.7}
及 \SGARefLink{par:sga3-VIA-5.3.1}{5.3.1} 的陈述可转述如下.
设 \(F\) 和 \(G\) 是在 \(A\) 上局部有限型且平坦的群,
其中 \(F\) 为 \(G\) 的闭正规子群. 映射
\[
  H\longmapsto F\backslash H,
  \qquad
  H'\longmapsto H'\times_{(F\backslash G)}G
\]
在 \(G\) 中包含 \(F\) 的平坦 \(A\)-子群与 \(F\backslash G\) 的平坦
\(A\)-子群之间定义了一个双射.
在这一双射下, \(G\) 中包含 \(F\) 的闭子群 (相应地, 正规子群)
对应于 \(F\backslash G\) 的闭子群 (相应地, 正规子群).%
\footnote{\label{note:sga3-VIA-52}编者注: 除了所引的第~IV~篇的陈述外,
还用到如下事实: 由于 \(G\to F\backslash G\) 忠实平坦,
\(G\) 的 \(A\)-子群 \(H\) 在 \(A\) 上平坦,
当且仅当 \(F\backslash H\) 在其上平坦.}

\medskip
\phantomsection\label{par:sga3-VIA-5.3.2}%
\noindent\textbf{5.3.2.---}
\SGARefLink{sga3:iv:proposition:5:2:9}{第~IV~篇命题~5.2.9} 给出如下结果.
设 \(F,H,G\) 是在 \(A\) 上局部有限型且平坦的群, 并设
\[
  F\subset H\subset G,
\]
其中 \(F\) 在 \(G\) 中闭, 在 \(H\) 中正规.
则 \(F\backslash H\) 在 \(F\backslash G\) 上自由地左作用, 商概形
\[
  (F\backslash H)\backslash(F\backslash G)
\]
存在, 且存在典范的 \(G\)-概形同构
\[
  (F\backslash H)\backslash(F\backslash G)
    \xrightarrow{\ \sim\ }
  H\backslash G.
\]

\medskip
\phantomsection\label{par:sga3-VIA-5.3.3}%
\noindent\textbf{5.3.3.---}
最后, \SGARefLink{sga3:iv:proposition:5:2:8}{第~IV~篇 5.2.8} 给出如下断言.
设 \(F,H,G\) 是在 \(A\) 上局部有限型且平坦的群.
假设 \(F\) 包含于 \(G\), 且在 \(G\) 中闭并正规,
\(H\) 包含于 \(G\), 而 \(F\cap H\) 在 \(A\) 上平坦. 以
\[
  F\mathbin{\overset{\tau}{\times}}_A H
\]
表示如下 \(A\)-群: 其底概形为 \(F\times_AH\), 乘法定义为
\[
  \bigl((x,h),(y,h')\bigr)
    \longmapsto
  \bigl(xhyh^{-1},hh'\bigr).
\]
同样, 设
\[
  u:H\cap F\longrightarrow
    F\mathbin{\overset{\tau}{\times}}_A H
\]
为单态射 \(x\mapsto(x^{-1},x)\), 并设 \(F\cdot H\) 为商
\[
  (F\cap H)\backslash
    \bigl(F\mathbin{\overset{\tau}{\times}}_A H\bigr).
\]
则存在典范同构
\[
  F\backslash(F\cdot H)
    \xrightarrow{\ \sim\ }
  (F\cap H)\backslash H.
\]
% Independent prose review: source and translation compared paragraph by paragraph through this component; statements, hypotheses, proofs, notes and diagram captions covered; bibliography reserved for its assigned reviewer; source math and reference keys unchanged.
